\documentclass[a4paper]{article}
\usepackage{amsfonts}
\usepackage{amsmath}
\usepackage{amssymb}
\usepackage{graphicx}     

\begin{document}
  
\noindent \large Auteur: Victoria Mazur \\
\noindent \large Studierichting: Informatica\\
\noindent \large Oefeningenreeks 6B nr. 5.7\\

\medskip

\normalsize

Bereken voor de volgende meetkundige rij $n$: \\

$x_3 = -36, x_n = -324, q = -3$ \\

$x_3 \cdot q^{n-3} = x_n$ \\

$\Leftrightarrow -36(-3)^{n-3} = -324$ \\

Beide zijden delen door $-36$: \\

$\Leftrightarrow (-3)^{n-3} = \dfrac{-324}{-36}$ \\

$\Leftrightarrow (-3)^{n-3} = 9$ \\

We weten dat $9 = 3^2 = (-3)^2$: \\

$\Leftrightarrow (-3)^{n-3} = (-3)^2$ \\

$\Rightarrow n - 3 = 2 \Leftrightarrow n = 5$ \\

\begin{thebibliography}{999}
%\bibitem{a} Referentie
\end{thebibliography}
\end{document}

